\section{Coloring}

\begin{definition}
    A \textbf{$k$-coloring} of $G = (V, E)$ is a mapping $c: V \to [k]$ such that $c(u) \neq c(v)$ for all $\{u,v\} \in E$.
\end{definition}

\begin{definition}
    The \textbf{chromatic number} $\chi(G)$ is the minimum number of colors.
\end{definition}

\begin{definition}
    A \textbf{stable set} is a set of vertices where no two vertices are adjacent.
\end{definition}

\begin{theorem}
    $\chi(G) \le 2 \iff G$ is bipartite.
\end{theorem}

\begin{theorem}
    For $k \ge 3$, the problem "Is $\chi(G) \le k$?" is NP-complete.
\end{theorem}

\begin{remark}
    For general graphs all $\epsilon$-approximations are NP-hard.
\end{remark}

\begin{algorithm}
#Greedy Coloring
#Requires at most maximum degree + 1 colors.

choose an arbitrary order of vertices v1...vn

c[v1] = 1
for i = 2 to n:
    c[vi] = min { j in [k] | j != c[vj] for all j < i and (vi, vj) in E }
\end{algorithm}

\begin{theorem}
    \textbf{Brooks' Theorem}: For a connected graph $G$ that is neither a complete graph $K_n$ nor an odd cycle $C_{2n+1}$:
    $$\chi(G) \le \Delta(G)$$
\end{theorem}

